\chapter{Conclusioni}
\label{cap:conclusioni}

\intro{Il presente capitolo tira le somme del lavoro svolto durante lo stage, confrontando la pianificazione con l'esecuzione reale, valutando il raggiungimento degli obiettivi, sintetizzando le conoscenze acquisite e proponendo possibili evoluzioni future.}\\

\section{Consuntivo finale}

Il progetto si è sviluppato sulle 8 settimane pianificate, per un totale di 300 ore. Il
monte ore è stato complessivamente rispettato, con una redistribuzione interna: le fasi di
studio (settimane 1-3) e di prototipazione (settimana 4) hanno richiesto un numero di ore
inferiore al preventivato grazie all'uso dell'agente AI come \emph{pair engineer}, permettendo così
un aumento del tempo a disposizione per lo sviluppo e l'approfondimento di funzionalità non
precedentemente preventivate:
\begin{itemize}
    \item \textbf{arricchimento automatico delle \emph{trace} con l'identità dell'utente
        autenticato}, ottenuto tramite un \emph{processor} \gls{otel} centralizzato in
        \texttt{BaseApplicationBuilder} che intercetta ogni \emph{span} radice e ne propaga
        l'\texttt{enduser.id} senza richiedere modifiche al codice dei singoli moduli;
    \item \textbf{integrazione del meccanismo degli \emph{exemplars}}, che collega ogni
        campione di metrica al \texttt{TraceId} della richiesta che lo ha generato,
        abilitando la navigazione diretta dalla dashboard Grafana al \emph{trace} di Tempo
        e completando la catena diagnostica
        metrica~$\rightarrow$~\emph{trace}~$\rightarrow$~utente~$\rightarrow$~\emph{log};
    \item \textbf{estensione dell'istrumentazione a moduli inizialmente non in scope}
        (\emph{Sia.Devextreme.Crud}, \emph{Sia.Devextreme.Kanban} e ulteriori moduli
        ancillari), realizzata tramite la \emph{Claude skill} \texttt{sia-telemetry-developer}
        che ha permesso la propagazione coerente del pattern \texttt{SiaXxxDiagnostics};
    \item \textbf{introduzione \emph{ex novo} di un'infrastruttura di testing} per il
        framework, comprensiva di \emph{unit test} sui \emph{controller}, \emph{integration
        test} di tipo \emph{smoke}, \emph{lifecycle} (CRUD/Grid) e \emph{summary}, con la
        relativa \emph{fixture} \texttt{SiaWebApplicationFactory} e i meccanismi di
        autenticazione/autorizzazione di test (\texttt{TestAuthHandler},
        \texttt{AllowAllPermissionCheckService}); la sola messa in funzione della prima
        \emph{suite} ha portato alla luce due difetti applicativi latenti;
    \item \textbf{popolamento delle dashboard Grafana con pannelli aggiuntivi} per la
        diagnosi a grana fine (latenza per endpoint, distribuzione degli errori per modulo,
        \emph{top users} per consumo di risorse), versionate nel repository tramite
        \emph{provisioning} dichiarativo;
    \item \textbf{produzione di una suite di \emph{Claude skill} procedurali} che
        codificano in forma riutilizzabile le convenzioni del framework (telemetria,
        \emph{controller}, \emph{service}, \emph{repository}, \emph{CRUD}, \emph{grid},
        ecc.), trasformando un esito incidentale dello stage in un asset metodologico
        permanente per il team.
\end{itemize}

% \section{Sintesi dei contributi originali}

% Prima di rendere conto del raggiungimento dei singoli obiettivi, è utile
% ricapitolare i contributi originali del lavoro, già anticipati nella
% sezione~\ref{sec:contributi}, evidenziandone l'impatto sul \emph{framework}
% \emph{Sia.Web}.

% Il primo contributo è la \textbf{progettazione e l'integrazione di un
% sistema di osservabilità completo} per un \emph{framework} che ne era
% sprovvisto. Il \emph{pipeline} \gls{otel} è stato centralizzato in
% \texttt{BaseApplicationBuilder} in modo tale che i moduli applicativi non
% debbano conoscerne i dettagli; ogni nuovo modulo aderisce alla telemetria
% semplicemente registrando il proprio \texttt{ActivitySource} e
% \texttt{Meter}. Questa scelta, unita alla tassonomia di attributi e metriche
% definita per i moduli core, riduce drasticamente l'attrito di adozione e
% garantisce confronti omogenei tra moduli e tra installazioni di clienti
% diversi.

% Il secondo contributo è l'introduzione di due meccanismi trasparenti al
% codice del modulo --- l'arricchimento automatico delle \emph{trace} con
% l'identità dell'utente autenticato e l'uso degli \emph{exemplar} ---, i
% quali abilitano una \textbf{catena diagnostica completa}
% metrica~$\rightarrow$~\emph{trace}~$\rightarrow$~utente~$\rightarrow$~\emph{log}.
% Senza questi meccanismi, l'osservabilità si ferma alla domanda ``\emph{quando
% e quanto} qualcosa è andato storto''; con essi diventa possibile rispondere
% a ``\emph{quale richiesta} è stata responsabile del picco di latenza, e
% \emph{a quale utente} stava rispondendo''.

% Il terzo contributo è la \textbf{nascita di un'infrastruttura di testing}
% nel \emph{framework}, introdotta da zero. La sola messa in funzione della
% prima \emph{suite} di \emph{integration test} ha portato alla luce due
% difetti applicativi latenti che si sarebbero manifestati esclusivamente
% presso un cliente specifico, in produzione: un risultato concreto che
% quantifica il valore della rete di sicurezza introdotta.

% Il quarto contributo, di natura metodologica, è la definizione di una
% \textbf{pratica di lavoro AI-assistita} disciplinata e replicabile: il
% \emph{workflow} \emph{Plan~$\rightarrow$~Confirm~$\rightarrow$~Implement},
% le \emph{skill} procedurali per la propagazione coerente di telemetria e
% test, l'adozione del \gls{mcp} per ridurre le ``allucinazioni'' lavorando
% su dati reali. Questa metodologia è stata utilizzata in modo sistematico
% durante lo stage e ha permesso di ottenere, in 300 ore, una copertura
% funzionale ed estensiva che difficilmente sarebbe stata raggiungibile con
% un approccio tradizionale, a parità di risorse.

\section{Raggiungimento degli obiettivi}

Gli obiettivi dichiarati nella sezione ``Obiettivi'' del capitolo~\ref{cap:descrizione-stage}
sono stati raggiunti:

\begin{itemize}
    \item l'analisi dei requisiti di osservabilità e la selezione delle tecnologie hanno
        portato all'adozione dello stack \gls{otel} + \emph{Grafana LGTM}, motivata dai
        criteri descritti nel capitolo~\ref{cap:tecnologie-adottate};
    \item l'architettura di telemetria è stata progettata e integrata con il framework
        \emph{.NET}, centralizzando la configurazione in \texttt{BaseApplicationBuilder}
        (capitolo~\ref{cap:progettazione});
    \item le strategie di campionamento sono state definite a due livelli (\emph{sampler}
        globale e \emph{processor} per-\emph{source}), con ratio configurabili per ambiente;
    \item gli strumenti di istrumentazione sono stati implementati per i moduli
        \emph{Sia.Auth}, \emph{Sia.Base} e \emph{Sia.Grid}, e propagati successivamente ad
        altri moduli del framework tramite le \emph{Claude skill} dedicate;
    \item la piattaforma di visualizzazione Grafana è stata configurata con i relativi
        \emph{backend} di \emph{storage} (Tempo, Loki, Prometheus) e con le dashboard
        versionate nel repository.
\end{itemize}

I requisiti funzionali, di qualità e di vincolo sono stati soddisfatti nella loro
totalità. Compresi quelli desiderabili e quelli aggiuntivi proposti dallo stagista stesso.
Consultabili nell'appendice~\ref{app:requisiti}.

\section{Conoscenze acquisite}

Lo stage ha portato all'acquisizione di competenze in più ambiti:

\begin{itemize}
    \item standard \gls{otel} e modello dei segnali (\emph{trace}, metriche, \emph{log},
        \emph{exemplars});
    \item stack \emph{Grafana} (Tempo, Prometheus, Loki), con i rispettivi linguaggi di
        query PromQL, TraceQL, LogQL;
    \item libreria Serilog e \emph{sink} OTLP, con correlazione \emph{log}-\emph{trace};
    \item applicazione dei \emph{design pattern} Adapter e Factory a un contesto reale di
        framework condiviso;
    \item metodologie di lavoro assistito da agenti AI, \emph{prompt} e \emph{skill
        engineering}, uso del \gls{mcp} per integrare l'agente con servizi runtime.
\end{itemize}

\section{Retention e gestione dei dati di telemetria}

Un aspetto operativo affrontato durante lo stage è la \emph{retention} dei dati di
telemetria. Ogni segnale ha caratteristiche di storage differenti:
\begin{itemize}
    \item le \emph{trace} su Tempo sono voluminose per singolo campione ma a bassa densità
        (sono campionate); la \emph{retention} suggerita è \emph{medio-breve} (7-14 giorni),
        sufficiente per la diagnosi di incidenti recenti;
    \item le metriche su Prometheus (o Mimir) sono aggregate e compatte: la \emph{retention}
        può estendersi a 90 giorni o più, permettendo il confronto di tendenze nel tempo;
    \item i \emph{log} su Loki sono voluminosi ma indicizzati per \emph{label}: la
        \emph{retention} tipica è di 30 giorni, con possibilità di archiviazione a lungo
        termine in storage a basso costo.
\end{itemize}
I tempi concreti di \emph{retention} sono configurati in funzione delle esigenze di indagine
e del costo di \emph{storage} accettabile per il singolo cliente.

\section{Valutazione personale e sviluppi futuri}

Dal punto di vista personale, lo stage ha permesso di lavorare a stretto contatto con un
\emph{framework} enterprise reale, cimentandosi sia con la progettazione (scelte di
cardinalità, \emph{sampling}, correlazione \emph{cross-signal}) sia con la
meta-programmazione del processo di sviluppo tramite \emph{Claude skill}. La collaborazione
con il \emph{team} aziendale e il confronto continuo sulle scelte architetturali sono stati
parte integrante del percorso formativo.

Tra gli sviluppi futuri più promettenti si collocano:
\begin{itemize}
    \item introduzione del \emph{tail-based sampling} tramite \emph{Collector}, che
        permetterebbe di conservare integralmente le trace di richieste lente o in errore
        anche con un \emph{ratio} globale basso;
    \item definizione di metriche \emph{RED} (Rate, Errors, Duration) e \emph{USE}
        (Utilization, Saturation, Errors) in forma automatizzata per ogni nuovo modulo;
    \item definizione di \emph{Service Level Objective} (SLO) e \emph{alerting} basato sul
        \emph{burn-rate} dell'\emph{error budget};
    \item integrazione di \emph{continuous profiling} tramite \emph{Pyroscope}, già presente
        nello stack Grafana, per correlare trace e profili di CPU/memoria;
    \item estensione del pattern di \emph{skill engineering} ad altri ambiti del framework,
        oltre alla telemetria (es.\ generazione di \emph{CRUD}, di \emph{controller}, di
        \emph{dashboard} di \emph{business}).
\end{itemize}
